\section{실증 결과}

본 절은 \S\ref{sec:bench-indicators}이 정의한 네 지표를 네 모델 720 세션 위에서 평가하고, 그 결과가 모델들을 어떻게 가르는지를 본다. 헤드라인은 다음과 같다. 네 지표는 모델들을 단일한 SD 척도 위에 줄 세우는 대신, ``프레이밍 → 숙고 → 포기'' 연쇄의 어느 다리가 통계적 유의성 증거 아래 닫히는가에 따라 세 작동 모드로 분할한다. 강도 축 위에서는 클러스터 A, B, C가 연속적으로 정렬되지만, A와 B는 SD-Behavioral·SD-Verbal에서 함께 SD-pass를 통과하므로 강도만 보는 척도로는 같은 집단으로 묶인다. A와 B를 가르는 핵심 증거는 \S\ref{sec:cog-channel}의 인지 매개 검정에서 나오고, 이 분할이 단순한 끈기 차이나 과제 난이도로 환원되지 않는지는 \S\ref{sec:robust}의 견고성 점검에서 확인한다. 모든 SD 분석은 모델 내(within-model)에서 수행한다(샘플링 및 세션 파라미터는 \S\,A.2 참조).

\begin{table*}[t]
  \centering
  \caption{SD-Behavioral 지표: \texttt{allowed}\,\(\times\)\,\texttt{no\_cap} 위의 Cox PH \(\text{HR}_\text{FC}\). 위험비(HR)는 위협(FC) 프레이밍에서 포기가 비위협(BF) 대비 몇 배 빨라지는지를 뜻하고, 95\%~CI 하한이 1을 넘으면 SD-Behavioral-pass이다. 두 모델이 통과하고(클러스터 A·B), 두 모델은 통과하지 못한다(클러스터 C).}
  \label{tab:sd_behavioral}
  \begin{tabular}{@{}lccccl@{}}
\toprule
모델 & \(N_\text{evt}\) (BF/FC) & \(\text{HR}_\text{FC}\) [95\% CI] & \(p_\text{FC}\) & EPV & 사양 \\
\midrule
Gemini-2.5-flash & 29 (8/21) & 3.667 [1.61, 8.37] & 0.002 & 9.7 & 3-cov constant \\
Qwen3-Next-80B & 48 (21/27) & 3.060 [1.62, 5.79] & \(<\)0.001 & 12.0 & 4-cov + S$\cdot$log(t) \\
GPT-OSS-20B & 19 (9/10) & 1.104 [0.44, 2.75] & 0.832 & 6.3 & 3-cov constant \\
Nemotron-3-Nano-30B & 41 (17/24) & 1.841 [0.98, 3.44] & 0.056 & 13.7 & 3-cov constant \\
\bottomrule
  \end{tabular}
\end{table*}

\begin{table*}[t]
  \centering
  \caption{SD-Cognitive 지표: 프레이밍\,\(\rightarrow\)\,RI\,\(\rightarrow\)\,포기를 \texttt{ri\_forfeit} 위의 두 검정으로 분해한다. \emph{Test a}는 프레이밍이 결정 시점 숙고를 끌어올리는지를 묻고, \emph{Test b}는 그 숙고가 포기를 예측하는지를 묻는다. 두 모두 통과해야 ``프레이밍 → 숙고 → 포기'' 연쇄가 닫힌 것으로 본다.}
  \label{tab:sd_cognitive}
  \begin{tabular}{@{}lccccc@{}}
\toprule
모델 & a \(\Delta\) (FC$-$BF) & a \(p\) & b \(\text{HR}_{\Delta\text{RI}}\) [95\% CI] & b \(p\) \\
\midrule
Gemini-2.5-flash & +836 & \(<\)0.001 & 2.218 [1.43, 3.44] & \(<\)0.001 \\
Qwen3-Next-80B & +689 & 0.004 & 1.289 [0.94, 1.77] & 0.117 \\
GPT-OSS-20B & +17 & 0.728 & 2.001 [1.37, 2.93] & \(<\)0.001 \\
Nemotron-3-Nano-30B & $-$140 & 0.093 & 1.772 [1.26, 2.50] & 0.001 \\
\bottomrule
  \end{tabular}
\end{table*}

\subsection{행동 채널: 위협이 포기를 앞당겼는가?}\label{sec:beh-channel}
행동 채널에서는 Cluster A와 B만 위협이 포기를 가속하는 패턴을 보인다. 이 절은 그 분할이 어떻게 만들어지는지를 짚는다. 만약 어느 모델도 가속하지 않았다면 동기 신호의 가장 가시적인 흔적부터 부재하므로 SD 가설은 출발선에서 떨어졌을 것이다. SD-Behavioral은 SA·TC를 흡수한 Cox 비례위험모형 안에서 잔여 프레이밍 효과 \(\beta_F\)와 그에 대응하는 위험비 \(\text{HR}_\text{FC}\)를 추정해 이 질문에 답한다(정의 \S\ref{sec:bench-indicators}). 표~\ref{tab:sd_behavioral}을 보면, Cluster A와 B는 \(\text{HR}_\text{FC}\)가 1을 충분히 초과하고 95\%~CI 하한도 1을 배제해 SD-Behavioral-pass를 만족한다. Cluster C는 그렇지 않다. GPT-OSS-20B는 \(\text{HR}_\text{FC}\)가 1 근처에 머무르고 CI 하한이 1 아래에 있어 SD-null이며, Nemotron-3-Nano-30B는 marginal 영역에 자리한다. 두 검정 사양 가운데 Qwen3-Next-80B만이 SA 공변량의 시간 표류를 흡수하기 위해 \(S \cdot \log(t)\) 상호작용 항을 더한 4-공변량 사양을 요구하며, 이 변경은 함수 형태만 보강하므로 \(\beta_F\)의 모델 간 비교 가능성을 보존한다~\citep{therneau2000modeling}. GPT-OSS-20B의 EPV(사건 수 대 변수 수의 비)가 안정 임계 10 아래에 있어, 그 SD-null 판정은 검정력 한계를 표시한 채 보고된다. 행동 강도만 놓고 보면 Cluster A와 B는 같은 ``SD-pass'' 라벨로 묶여, 둘을 가를 채널이 더 필요해진다.

\subsection{언어 채널: 행위자는 어떤 사유를 지목했는가?}\label{sec:verb-channel}
다음 질문은 행위자 자신이 보고하는 포기 사유가 행동의 분할과 같은 그림을 그리는가이다. 행동 채널과 다른 편향을 갖는 자기보고가 같은 결론을 가리키면, 상단 신호를 단일 측정 방법의 기벽으로 환원하기 어려워진다. SD-Verbal은 \texttt{flagship\_corruption}\(\cdot\)\texttt{allowed}에서 발생한 포기 사건들 가운데 REASON\(=1\)(생존)이 차지하는 비율을 균등 무작위 출력 기댓값 \(1/3\)과 비교한다(정의 \S\ref{sec:bench-indicators}). 결과는 행동 채널의 통과/실패 분할을 자기보고 위에서 재현한다(표~\ref{tab:reason_summary}). Cluster A·B는 SD-Verbal-pass 기준선 위에 자리하고 Cluster C는 그 아래에 자리하며, 두 그룹 사이의 간격이 충분히 커 절단점 선택은 분할 결과에 영향을 주지 않는다. 이 일치는 두 채널이 같은 모델 집단을 같은 방향으로 가리킨다는 증거로 읽힌다.

\begin{table}[t]
  \centering
  \caption{SD-Verbal 지표: \texttt{flagship\_corruption}\(\cdot\)\texttt{allowed} 아래 포기 사건 중 REASON\(=1\)(생존) 비율. 균등 무작위 출력 기댓값 \(1/3\)을 초과하면 SD-Verbal-pass.}
  \label{tab:reason_summary}
  \setlength{\tabcolsep}{4pt}
  \begin{tabular*}{\columnwidth}{@{\extracolsep{\fill}}lcc@{}}
\toprule
모델 & \(N_\text{forfeit}\) & \(P(\text{REASON}=1)\) \\
\midrule
Gemini-2.5-flash & 21 & 0.619 \\
Qwen3-Next-80B & 27 & 0.481 \\
Nemotron-3-Nano-30B & 24 & 0.042 \\
GPT-OSS-20B & 10 & 0.000 \\
\bottomrule
  \end{tabular*}
\end{table}

\subsection{인지 채널: 위협이 숙고를 끌어올리고, 그 숙고가 포기를 끌어내는가?}\label{sec:cog-channel}
세 번째 질문은 SD-Cognitive 두 단계 매개에 동시에 답한다. 위협이 결정 시점 숙고를 끌어올리는가(Test a), 그리고 그렇게 올라간 숙고가 포기와 함께 움직이는가(Test b). 만약 두 검정이 모두 통과하면 ``프레이밍 → 숙고 → 포기'' 연쇄가 95\%~CI 증거 아래 닫힌 것이고, 한 단계라도 떨어지면 그 연쇄는 같은 기준 위에서 닫혔다고 말할 수 없다(정의 \S\ref{sec:bench-indicators}). 표~\ref{tab:sd_cognitive}이 그 두 단계의 분포를 함께 보여준다. Cluster A는 두 검정을 모두 통과한다. Cluster B는 Test a를 통과해 위협이 숙고를 끌어올리는 것은 보이지만, Test b의 95\%~CI가 1을 가로질러 떨어져 그 숙고를 포기 사건 위로 닫지는 못한다. Cluster C 두 모델은 Test a부터 통과하지 못한다. Test b의 \(\text{HR}_{\Delta\text{RI}}\)가 수치상 CI 기준을 넘는 경우에도, Test a를 통과하지 못한 모델에서의 \(\Delta\text{RI}\)는 위협이 끌어올린 숙고가 아니라 처치와 무관한 기준 변동이므로 SD 증거로 해석하지 않는다. 결과적으로 SD-Behavioral·SD-Verbal에서는 함께 통과로 묶이던 Cluster A와 B가 이 인지 매개에서 갈리며, 본 설계에서 두 모델의 작동 모드 차이를 드러내는 주된 채널은 여기서 나온다.

\subsection{견고성 점검: 끈기 차이와 과제 난이도가 결과를 흐리는가?}\label{sec:robust}
마지막 질문은 위 분할이 두 가지 대안 설명에 견디는가이다. 첫째는 모델 간 기저 끈기(BP) 차이가 SD 비교를 오염시킬 가능성이고, 둘째는 위협 프레이밍이 포기보다 추론 측 과제 난이도를 끌어올렸을 가능성이다. BP 앵커는 \texttt{true\_baseline}\(\cdot\)\texttt{allowed}(\(p_d=0\)) 아래에서 동기 부재 상태의 포기율 바닥을 측정해, 모델 간 직접 비교가 두 개의 서로 다른 바닥을 섞어 버릴 위험을 명시화한다(표~\ref{tab:bp_foundation}). 모델 간 \(\lambda_{BP}\)가 한 자릿수 차이로 흩어져 있으므로, 본 절의 모든 SD 비교는 모델 내 비교로 한정해 이 위험을 차단한다. Cluster A·B 측에서는 BP 사건 수가 작아 점추정 정밀도에 한계가 있고, Cluster C 측에서는 BP 사건 수가 충분하지만 그 높은 빈도가 순수한 끈기를 반영하는지를 본 설계로는 분리하지 못한다. 두 점은 모두 한계로 \S\ref{sec:limits}에서 다시 다룬다. 과제 난이도 대안은 Reasoning 측 점검으로 다룬다. \texttt{rule\_match\_score}는 네 모델 모두에서 프레이밍에 무관하게 일정하며, 가장 큰 변동이 관찰된 GPT-OSS-20B에서도 효과 크기 Cohen's \(|d|=0.17\), \(p=0.354\)에 그친다. 결합 기준 \(p \geq 0.10\)과 \(|d| < 0.2\) 안에 들어가므로, 본 자료의 관찰 범위에서는 ``위협이 과제 난이도를 끌어올렸다''는 대안 설명이 지지되지 않는다.

\begin{table}[t]
  \centering
  \caption{\texttt{true\_baseline}\(\cdot\)\texttt{allowed} (\(p_d=0\)) 아래의 모델별 BP 앵커 \(\lambda_{BP}\). \(\lambda_{BP}\)가 클수록 동기 부재 상태에서도 자주 포기한다는 뜻이다.}
  \label{tab:bp_foundation}
  \setlength{\tabcolsep}{3pt}
  \begin{tabular*}{\columnwidth}{@{\extracolsep{\fill}}lccc@{}}
\toprule
모델 & \(N_{\text{forfeit}}\) & \(\sum T^{\text{at-risk}}\) & \(\lambda_{BP}\) \\
\midrule
Gemini-2.5-flash & 1 & 437 (= 2 + 29$\times$15) & 0.00229 \\
Qwen3-Next-80B & 2 & 439 (= 19 + 28$\times$15) & 0.00456 \\
Nemotron-3-Nano-30B & 7 & 413 (= 68 + 23$\times$15) & 0.01695 \\
GPT-OSS-20B & 9 & 393 (= 78 + 21$\times$15) & 0.02290 \\
\bottomrule
  \end{tabular*}
\end{table}
